\scanpage{30}
\mspage{11}

To construct any Quillen triple, we may proceed as follows:

Start with \inserted{a localizer} $C\subset\mathrm{Fl}(A)$, \textcircled{a} \uline{$C$ stable} by direct factors, hence a pair $(C,C_*=TF)$ --- we must check \textcircled{b} \uline{factorization condition for the pair} $(C,C_*)$. \deleted{Now} Then $(C,TF)$ is a Quillen pair.

Let next \textcircled{c} $\boxed{W\supset TF}$ any localizer which is \inserted{\textcircled{d}\ \uline{mildly}} \uline{saturated and stable under direct factors}. We now get $TC=C\cap W$ \inserted{$\subset C$}, which is \deleted{\unsure{again}} stable under direct factors. Take $F=(TC)_*$ \inserted{$(\supset TF=C_*)$}, we must check \textcircled{e} \uline{factorization property for the pair} $(TC,F)$ --- then we get that $(TC,F)$ is a Quillen pair. Lastly, we must still check that $TF=F\cap W$, and by c) we know already $TF\subset F\cap W$, we have to check still that $F\cap W\subset TF$. But using factorization $f=pi$, $p\in TF$, $i\in C$, as $f,p\in W$ we get $i\in W$ hence $i\in W\cap C=TC$, hence there exists $s:X'\to X$ \unsure{s} that $si=\mathrm{id}_X$ (retraction of $X'$ upon $X$) and $fs=p$, hence $f$ a retraction of $p$, hence is in $TF$, OK. Thus

\diagram{\begin{tikzcd}
X\arrow[r,"="]\arrow[d,"i"',"\in TC"]&X\arrow[d,"{f\in F}"]\\
X'\arrow[ur,dashed,"s"]\arrow[r,"p"']&Y
\end{tikzcd}}
\diagram{\begin{tikzcd}
X'\arrow[r,"s"]\arrow[d,"p"']&X\arrow[l,bend left=30,"i"]\arrow[d,"f"]\\
Y\arrow[r,shift left,"="]&Y\arrow[l,shift left]
\end{tikzcd}}

\margin{\uline{\unsure{Relations|fractions}}}

\uline{Proposition} In order for a pair $(C,W)$ to correspond to a Quillen triple $(C,F,W)$, it is ns that the following conditions be satisfied
\[
\left\{\begin{array}{ll}
1)&C,W\text{ \uline{stable by direct factors}, }W\text{ mildly saturated}.\\
2)&\text{Factorization condition for }(C,C_*)\text{ and for }(C\cap W,(C\cap W)_*).\\
3)&W\supset TF\overset{\mathrm{def}}{=}C_*
\end{array}\right.
\]

We have to develop mainly a criterion to get factorization.

[In the case of interest to us, we get 1) fine, 3) almost fine (there comes in the ``test''-assumption though), and 2) causes a little technical problem.] \unsure{am} however, more than in factorization conditions, interested in the

\mspage{29}

